\documentclass{article}
\usepackage{amsmath, amssymb, amsthm}
\usepackage{hyperref}
\usepackage{geometry}
\geometry{margin=1in}

\title{Erd\H{o}s Problem \#293}
\date{Last updated: 2025-08-31}
\author{Source: erdosproblems.com}

\begin{document}
\maketitle

\noindent\textbf{Status:} OPEN \quad \textbf{No prize}

\noindent\textbf{Tags:} number theory, unit fractions

\medskip

\noindent\textbf{Problem Statement:}

Let $k\geq 1$ and let $v(k)$ be the minimal integer which does not appear as some $n_i$ in a solution to\[1=\frac{1}{n_1}+\cdots+\frac{1}{n_k}\]with $1\leq n_1&#60;\cdots &#60;n_k$. Estimate the growth of $v(k)$.




Results of Bleicher and Erd\H{o}s \cite{BlEr75} imply $v(k) \gg k!$. It may be that $v(k)$ grows doubly exponentially in $\sqrt{k}$ or even $k$. An elementary inductive argument shows that $n_k\leq ku_k$ where $u_1=1$ and $u_{i+1}=u_i(u_i+1)$, and hence\[v(k) \leq kc_0^{2^k},\]where\[c_0=\lim_n u_n^{1/2^n}=1.26408\cdots\]is the 'Vardi constant' (small improvements on this are possible as in [148] ). van Doorn and Tang \cite{vDTa25b} have proved that\[v(k)\geq e^{ck^2}\]for some constant $c&gt;0$, and noted a close connection to [304] . In particular, if $N(b)\ll \log\log b$ as in [304] then it is likely the methods of \cite{vDTa25b} prove $v(k) \geq e^{e^{ck}}$ for some $c&gt;0$.




References

[BlEr75] Bleicher, M. N. and Erd\H{o}s, P., The number of distinct subsums of $\sum \sb{1}\spN\,1/i$ . Math. Comp. (1975), 29-42.

[vDTa25b] W. van Doorn and Q. Tang, The smallest denominator not contained in a unit fraction decomposition of 1 with fixed length . arXiv:2512.22083 (2025).

\noindent\textbf{Related OEIS sequences:} possible


\bigskip
\noindent\small{Source: \url{https://www.erdosproblems.com/293}}
\end{document}
